\documentclass[12pt,a4paper]{article}
\usepackage[margin=25mm, headheight=15pt, headsep=10pt]{geometry}
\usepackage{fontspec}
\usepackage[table]{xcolor} % Note: table option is required for rowcolors
\usepackage{titlesec}
\usepackage{tocloft}
\usepackage{fancyhdr}
\usepackage{booktabs}
\usepackage{tcolorbox}
\tcbuselibrary{skins, breakable}
\usepackage[colorlinks=true, linkcolor=emerald, urlcolor=emerald, bookmarks=true]{hyperref}

\definecolor{emerald}{RGB}{0,102,51}
\definecolor{darkred}{RGB}{139,0,0}
\definecolor{lightgray}{RGB}{245,245,245}

\usepackage[nil]{babel}
\babelprovide[import, main]{bengali}
\babelprovide[import]{arabic}
\babelfont{rm}[Renderer=HarfBuzz,Scale=1.0]{Noto Serif Bengali}
\babelfont{sf}[Renderer=HarfBuzz]{Noto Sans Bengali}
\babelfont{tt}[Renderer=HarfBuzz]{Noto Sans Bengali}
\babelfont[arabic]{rm}[Renderer=HarfBuzz,Scale=1.1]{Noto Naskh Arabic}

% Section styling
\titleformat{\section}{\Large\bfseries\color{emerald}}{\thesection.}{0.5em}{}
\titleformat{\subsection}{\large\bfseries\color{emerald!85!black}}{\thesubsection.}{0.5em}{}
\titleformat{\subsubsection}{\normalsize\bfseries\color{emerald!70!black}}{\thesubsubsection.}{0.5em}{}
\titlespacing*{\section}{0pt}{18pt}{8pt}

% Fancy Header/Footer setup
\pagestyle{fancy}
\fancyhf{} % clear all header and footer fields
\fancyhead[L]{\color{emerald}\small\textbf{\nouppercase{\leftmark}}}
\fancyfoot[L]{\scriptsize\color{black!50}{\lat{AI}}-সংকলিত, আলেম-যাচাইকৃত নয়}
\fancyfoot[C]{\color{emerald!70!black}\thepage}
\renewcommand{\headrulewidth}{0.4pt}
\renewcommand{\headrule}{\hbox to\headwidth{\color{emerald}\leaders\hrule height \headrulewidth\hfill}}
\renewcommand{\footrulewidth}{0pt}

% Custom boxes
% 1. Ayat/Hadith Box
\newtcolorbox{ayatbox}[1][]{
    enhanced, breakable, 
    colback=emerald!5, 
    colframe=emerald, 
    boxrule=0.5pt, 
    arc=3pt, 
    left=10pt, right=10pt, top=10pt, bottom=10pt,
    #1
}

% 2. Quote/Translation Box
\newtcolorbox{quotebox}[1][]{
    enhanced, breakable,
    frame hidden,
    colback=lightgray,
    borderline west={3pt}{0pt}{emerald},
    left=15pt, right=10pt, top=10pt, bottom=10pt,
    sharp corners,
    #1
}

% 3. AI-generated notice box
\newtcolorbox{warnbox}[1][]{
    enhanced, breakable,
    colback=red!4,
    colframe=darkred,
    boxrule=0.8pt,
    arc=3pt,
    left=10pt, right=10pt, top=10pt, bottom=10pt,
    #1
}

% Commands
\newcommand{\arab}[1]{\foreignlanguage{arabic}{#1}}
\newcommand{\kib}[1]{\textcolor{emerald}{\textbf{#1}}}

% Latin (English) fallback: Noto Serif Bengali has NO Latin glyphs
\newfontfamily\latfont[Renderer=HarfBuzz]{Noto Serif}
\newcommand{\lat}[1]{{\latfont #1}}

\setlength{\parskip}{5pt}
\setlength{\parindent}{0pt}

% Fix for missing bullet in Noto Serif Bengali
\renewcommand{\labelitemi}{$\bullet$}



\begin{document}
\begin{center}{\Large\bfseries\color{emerald} পার্ট ২ — নামাজের পূর্বশর্ত: পবিত্রতা, সময়, কিবলা, পোশাক, আযান ও নিয়ত}\end{center}
\vspace{1em}
\section*{পার্ট ২ — নামাজের পূর্বশর্ত: পবিত্রতা, সময়, কিবলা, পোশাক, আযান ও নিয়ত}

\begin{warnbox}
 \textbf{গুরুত্বপূর্ণ নোটিশ (\lat{Important} \lat{Notice}):} এই কন্টেন্টটি \lat{AI} (কৃত্রিম বুদ্ধিমত্তা) দিয়ে প্রস্তুত ও সংকলিত। \lat{AI} কঠোর সূত্র-মান নীতি মেনে শুধু নির্ভরযোগ্য উৎস থেকে তথ্য সংগ্রহ করেছে — কেবল সহীহ/হাসান হাদিস ও সহীহ সনদের আসার রাখা হয়েছে, দুর্বল সনদ বাদ দেওয়া হয়েছে। তবে এটি কোনো আলেম বা মুহাদ্দিস কর্তৃক যাচাইকৃত নয়।
\end{warnbox}

\begin{quote}
\textbf{পড়ার নিয়ম:} নামাজ শুরু করার আগে যা যা \textbf{নিশ্চিতভাবে} ঠিক থাকতে হবে — সেগুলোই এই অংশের বিষয়। প্রতিটি শর্তের সাথে থাকবে \textbf{কুরআনের আয়াত / সহীহ হাদিস} $\rightarrow$ \textbf{মাযহাবগত অবস্থান (যেখানে আছে)} $\rightarrow$ \textbf{ব্যবহারিক নিয়ম}। \\
\textbf{সম্পর্কিত ফাইল:} পার্ট ৩ (ফরজ-ওয়াজিব-সুন্নতের শ্রেণিবিভাগ), পার্ট ৪ (পূর্ণ পদ্ধতি), পার্ট ৫গ (সফর-অসুস্থ-নারী), পার্ট ৭ক (নামাজ ভঙ্গকারী বিষয়)।
\end{quote}

\medskip\hrule\medskip

\subsection*{১. শর্ত (\lat{Conditions}) কী — এক নজরে}

নামাজের \textbf{বৈধতা} নির্ভর করে কিছু পূর্বশর্তের উপর। এগুলো থাকলে নামাজ শুরুর যোগ্যতা হয়; এর কোনোটি না থাকলে নামাজ শুরুই হয় না বা শুরু হলেও গ্রহণযোগ্য হয় না। ফকিহরা শর্তগুলো এভাবে তালিকাবদ্ধ করেন (সব মাযহাবে একই ক্রম নয়, তবে মূল বিষয়গুলো এক):

\begin{table}[h]\centering\small
\begin{tabular}{lll}
\# & শর্ত & ভিত্তি \\
\hline
১ & ইসলাম & কুরআন-সুন্নাহয় ইবাদত কেবল মুমিনের উপর ফরজ \\
২ & সুস্থ বুদ্ধি (আকল) & \lat{“}তিন ব্যক্তি থেকে কলম উঠিয়ে নেওয়া হয়েছে…\lat{”} (তিরমিজি ১৪২৩ — হাসান) \\
৩ & বয়স — বালেগ হওয়া (তাময়ীজ) & একই হাদিস; নাবালেগের উপর ফরজ নয়, তবে শিক্ষার জন্য আদেশযোগ্য \\
৪ & পবিত্রতা — ছোট-বড় হাদাস থেকে & সূরা মায়িদাহ ৫:৬; সহীহ বুখারি ১৩৫ \\
৫ & শরীর, কাপড় ও স্থান নাপাক (নাজাস) থেকে পবিত্র & সূরা আল-মুদাসসির ৭৪:৪-৫; সহীহ মুসলিম ৫৫৩ (মতভেদসহ) \\
৬ & সতর (আওরাত) ঢাকা & সূরা আল-আ'রাফ ৭:৩১; সুনানে আবু দাউদ ৬৪১ (সহীহ) \\
৭ & কিবলামুখী হওয়া & সূরা আল-বাকারাহ ২:১৪৪ \\
৮ & নির্দিষ্ট সময় (ওয়াক্ত) & সূরা আন-নিসা ৪:১০৩; তিরমিজি ১৪৯ (হাসান) \\
৯ & নিয়ত & সহীহ বুখারি ১ \\
\hline
\end{tabular}
\end{table}

\begin{quote}
 \textbf{গুরুত্বপূর্ণ:} উপরের ৯টি শর্তের প্রতিটির গুরুত্ব ও ভিত্তি ভিন্ন। যেমন — \textbf{ইসলাম ও সুস্থ বুদ্ধি} এমন মৌলিক শর্ত যাতে কেউ ব্যর্থ হলে নামাজের কথা-ই ওঠে না; আবার \textbf{নিয়ত} সম্পর্কে ফকিহদের মধ্যে অবস্থান ভিন্ন (কারো মতে শর্ত, কারো মতে রুকন — পার্ট ৩ দেখুন)। তাই কোনো শর্তকে 'অনুচ্ছেদগতভাবে সমান' ভাবা ঠিক হবে না।
\end{quote}

\textbf{শর্ত বনাম রুকন — একটি স্পষ্ট পার্থক্য:}

\begin{itemize}
\item \textbf{শর্ত (\arab{شرط}):} নামাজের \textbf{আগে} থাকতে হয় — যেমন ওযু, সময়, কিবলা, সতর। শর্ত না থাকলে নামাজ শুরুই হয় না।
\item \textbf{রুকন (\arab{رکن}):} নামাজের \textbf{মাঝে} সম্পাদন করতে হয় — যেমন কিয়াম, রুকু, সিজদা, তাশাহহুদ। রুকন বাদ গেলে নামাজ ভেঙে যায়।
\end{itemize}
\medskip\hrule\medskip

\subsection*{২. পবিত্রতা (তাহারাত) — সবচেয়ে গুরুত্বপূর্ণ পূর্বশর্ত}

পবিত্রতা দুই প্রকার, দুইটিই নামাজের জন্য আবশ্যক:

\begin{enumerate}
\item \textbf{হাদাস (\arab{حدث}):} অদৃশ্য, বিধানগত অপবিত্রতা — যা ওযু/গোসল/তায়াম্মুমের মাধ্যমে দূর হয়। দুই ভাগ:
\end{enumerate}
\begin{itemize}
\item \textbf{হাদাস আসগর (ছোট):} ওযু ভঙ্গকারী অবস্থা — ওযু দূর করে।
\item \textbf{হাদাস আকবার (বড়):} জুনুবি, হায়েয, নিফাস — গোসল দূর করে।
\end{itemize}
\begin{enumerate}
\item \textbf{নাজাস (\arab{نجس}):} দৃশ্যমান অপবিত্রতা — যেমন প্রস্রাব, মল, রক্ত। এটি শরীর, কাপড় ও স্থান থেকে ধুয়ে দূর করতে হয়।
\end{enumerate}
\subsubsection*{২.১ ওযু}

\textbf{মূল আয়াত (সূরা আল-মায়িদাহ ৫:৬):}

\begin{center}\arab{يَـٰٓأَيُّهَا ٱلَّذِينَ ءَامَنُوٓا۟ إِذَا قُمْتُمْ إِلَى ٱلصَّلَوٰةِ فَٱغْسِلُوا۟ وُجُوهَكُمْ وَأَيْدِيَكُمْ إِلَى ٱلْمَرَافِقِ وَٱمْسَحُوا۟ بِرُءُوسِكُمْ وَأَرْجُلَكُمْ إِلَى ٱلْكَعْبَيْنِ ۚ وَإِن كُنتُمْ جُنُبًا فَٱطَّهَّرُوا۟ ۚ وَإِن كُنتُم مَّرْضَىٰٓ أَوْ عَلَىٰ سَفَرٍ أَوْ جَآءَ أَحَدٌ مِّنكُم مِّنَ ٱلْغَآئِطِ أَوْ لَـٰمَسْتُمُ ٱلنِّسَآءَ فَلَمْ تَجِدُوا۟ مَآءً فَتَيَمَّمُوا۟ صَعِيدًا طَيِّبًا فَٱمْسَحُوا۟ بِوُجُوهِكُمْ وَأَيْدِيكُم مِّنْهُ ۚ مَا يُرِيدُ ٱللَّهُ لِيَجْعَلَ عَلَيْكُم مِّنْ حَرَجٍ وَلَـٰكِن يُرِيدُ لِيُطَهِّرَكُمْ وَلِيُتِمَّ نِعْمَتَهُۥ عَلَيْكُمْ لَعَلَّكُمْ تَشْكُرُونَ}\end{center}

\textbf{উচ্চারণ:} ইয়া- আয়্যুহাল্লাযীনা আ-মানূ, ইযা- কুমতুম ইলাস-সলাতি ফাগছিলূ উজূহাকুম ওয়া আইদিয়াকুম ইলাল-মারাফিক্বি, ওয়ামসাহূ বিরুঊসিকুম, ওয়া আরজুলাকুম ইলাল-কা'বাইন। ওয়া ইন কুনতুম জুনুবান ফাত্বাহহারূ। ওয়া ইন কুনতুম মারদ্বা- আও 'আলা- সাফারিন, আও জাআ আহাদুম মিনকুম মিনাল-গাইত্বি, আও লা-মাসতুমুন-নিসাআ, ফালাম তাজিদূ মাআন, ফাতায়াম্মামূ স'য়ীদান ত্বাইয়িবান, ফামসাহূ বিউজূহিকুম ওয়া আইদীকুম মিনহু। মা- ইউরীদুল্লা-হু লি-ইয়াজ'আলা 'আলাইকুম মিন হারাজিন, ওয়া-লাকিন ইউরীদু লি-ইউত্বাহহিরাকুম ওয়া লি-ইউতিম্মা নি'মাতাহূ 'আলাইকুম লা'আল্লাকুম তাশকুরূন।

\textbf{বাংলা অর্থ (\lat{pure}):}
\textbf{\lat{“}হে ঈমানদারগণ! যখন তোমরা নামাজের জন্য দাঁড়াতে চাও, তখন তোমাদের মুখমণ্ডল ও কনুই পর্যন্ত দুই হাত ধুয়ে নাও, মাথায় মাসেহ করো, এবং গোড়ালি পর্যন্ত দুই পা ধুয়ে নাও। আর যদি তোমরা জুনুবি (বড় অপবিত্র) হও, তবে গোসল করে পবিত্র হও। আর যদি তোমরা অসুস্থ থাকো, বা সফরে থাকো, বা তোমাদের কেউ পায়খানা থেকে এসে থাকে, অথবা তোমরা স্ত্রীদের সাথে সহবাস করে থাকো, আর (কোথাও) পানি না পাও, তবে পবিত্র মাটি দিয়ে তায়াম্মুম করো — তা দিয়ে তোমাদের মুখমণ্ডল ও দুই হাত মাসেহ করো। আল্লাহ তোমাদের উপর কোনো সংকীর্ণতা বানাতে চান না; বরং তিনি তোমাদের পবিত্র করতে চান এবং তোমাদের প্রতি তাঁর নেয়ামত পূর্ণ করতে চান, যাতে তোমরা শুকরিয়া আদায় করো।\lat{”}}

\textbf{সহীহ হাদিস (সহীহ বুখারি ১৩৫):}

\begin{quote}
আবু হুরাইরা (রাঃ) থেকে: রাসূলুল্লাহ (সা.) বলেছেন — \textbf{\lat{“}ওযু না হওয়া পর্যন্ত হাদাসগ্রস্ত ব্যক্তির নামাজ কবুল হয় না।\lat{”}} (হাদ্রামৌতের এক ব্যক্তি জিজ্ঞেস করলেন: হাদাস কী? আবু হুরাইরা বললেন: বায়ু বের হওয়া।)
\end{quote}

\#\#\#\# ওযুর ফরজ — চার মাযহাবের তুলনা

\begin{table}[h]\centering\small
\begin{tabular}{lllll}
কাজ & হানাফি & মালেকি & শাফেয়ি & হাম্বলি \\
\hline
নিয়ত & সুন্নাত & ফরজ & ফরজ & ফরজ \\
চেহারা ধোয়া & ফরজ & ফরজ & ফরজ & ফরজ \\
দুই হাত কনুইসহ & ফরজ & ফরজ & ফরজ & ফরজ \\
মাথা মাসেহ & ফরজ (চতুর্থাংশ) & ফরজ (পুরো মাথা) & ফরজ (সামান্য অংশ) & ফরজ (পুরো মাথা) \\
দুই পা গোড়ালিসহ & ফরজ & ফরজ & ফরজ & ফরজ \\
ক্রম (তরতীব) & সুন্নাত & ফরজ & ফরজ & ফরজ \\
ধারাবাহিকতা (মুয়ালাত) & সুন্নাত & ফরজ & সুন্নাত & ফরজ \\
ঘষে ঘষে ধোয়া (দালক) & — & ফরজ & সুন্নাত & সুন্নাত \\
\hline
\end{tabular}
\end{table}

\begin{quote}
\textbf{ব্যাখ্যা:} সব মাযহাবের মিল — \textbf{চারটি অঙ্গ ধোয়া/মাসেহ করা ফরজ} (চেহারা, দুই হাত কনুইসহ, মাথা, দুই পা)। পার্থক্য মূলত \textbf{মাথার মাসেহের পরিমাণ} ও \textbf{নিয়ত/ক্রম/ধারাবাহিকতার} প্রকৃতিতে। এগুলো আংশিক \textbf{শব্দগত-প্রায়} পার্থক্য, কারণ বেশিরভাগ ক্ষেত্রে বাস্তব অনুশীলন একই — তবে মাসেহের পরিমাণ ও নিয়তের হুকুমে \textbf{বাস্তব পার্থক্য} আছে।
\end{quote}

\#\#\#\# ওযু ভঙ্গকারী কারণ (নাকিস/মুবতিলাত)

\textbf{সর্বসম্মত (ইজমায়ি):}
\begin{enumerate}
\item সামনে বা পেছনের পথ থেকে কিছু বের হওয়া — প্রস্রাব, পায়খানা, বায়ু।
\item অচেতন হওয়া, পাগল হওয়া, মদ্যপানে নেশা — যে কারণে জ্ঞান লোপ পায়।
\item (গভীর) ঘুম — যে ঘুমে আসন-স্পর্শ আলগা হয়ে যায়। সামান্য ঢুলুনি নিয়ে মতভেদ আছে।
\end{enumerate}
\textbf{মাযহাবগত পার্থক্যসহ:}

\begin{table}[h]\centering\small
\begin{tabular}{lllll}
কারণ & হানাফি & মালেকি & শাফেয়ি & হাম্বলি \\
\hline
সরাসরি (বস্ত্রহীন) জননাঙ্গ স্পর্শ & ভাঙে & ভাঙে না & ভাঙে & ভাঙে \\
স্ত্রীকে কামভাবসহ (সরাসরি) স্পর্শ & ভাঙে & ভাঙে (কামভাব থাকলে) & ভাঙে না & ভাঙে না \\
শরীর থেকে প্রবাহিত রক্ত/পুঁজ/বমি & ভাঙে & ভাঙে (প্রবাহিত হলে) & ভাঙে না & ভাঙে না \\
নামাজে উচ্চস্বরে হাসি & ওযু ও নামাজ দুটিই ভাঙে & নামাজ ভাঙে, ওযু নয় & নামাজ ভাঙে, ওযু নয় & নামাজ ভাঙে, ওযু নয় \\
উটের গোশত খাওয়া & ভাঙে & ভাঙে না & ভাঙে না & ভাঙে \\
\hline
\end{tabular}
\end{table}

\begin{quote}
\textbf{কেন পার্থক্য?} — \lat{“}যে নিজের জননাঙ্গ স্পর্শ করবে সে যেন ওযু করে\lat{”} (সুনানে আবু দাউদ ১৮১, তিরমিজি ৮২ — সহীহ/হাসান) — এই হাদিস নিয়ে মালেকি মাযহাব মনে করে সেটি নির্দিষ্ট বাস্তবায়নে ছিল বা বিধান হিসেবে বহাল নয়; বাকি তিন মাযহাব হাদিসের জাহেরি (আপাত) অর্থ অনুযায়ী ওযু ভাঙার বিধান দিয়েছেন। আবার \textbf{রক্তের} ক্ষেত্রে হানাফি ও মালেকি \lat{“}শরীর থেকে বের হওয়া নাজিস\lat{”}-কে হাদাস ধরেছেন, শাফেয়ি ও হাম্বলি তা স্পর্শ-পবিত্রতা হিসেবে দেখেছেন। \textbf{উটের গোশত} সম্পর্কে সহীহ হাদিস আছে — \lat{“}উটের গোশত খেলে ওযু করো\lat{”} (সহীহ মুসলিম ৩৬০)। হানাফি ও হাম্বলি এই হাদিসের জাহেরি অর্থ নিয়েছেন; মালেকি ও শাফেয়ি মনে করেন এখানে 'ওযু' দ্বারা দাঁত-মুখ পরিষ্কার (নক-ভুবন) বোঝানো হয়েছে অথবা বিধানটি রহিত।
\end{quote}

\begin{quote}
\textbf{সতর্কতা:} কোনো একটি মাযহাবের অবস্থানকে 'একমাত্র সঠিক' বলা যাবে না — এগুলো বৈধ ইখতিলাফ। একজন মুকাল্লিদ (যে কোনো একটি মাযহাব মেনে চলে) তার মাযহাবের বিধানই আমল করবে; তবে আপতকালীন সন্দেহে নির্ভরযোগ্য আলেমের পরামর্শ নেবে।
\end{quote}

\#\#\#\# মাসেহ খুফফাইন (মোজার উপর মাসেহ)

মোজা বা খুফ (চামড়ার মোজা) পরার পর ওযু করলে, পরবর্তী ওযুতে \textbf{মোজার উপর মাসেহ} করা জায়েয — তবে পায়ের উপর নয়। হাদিস:

\begin{quote}
আল-মুগীরা ইবনে শুবা (রাঃ) থেকে: রাসূলুল্লাহ (সা.) ওযুর সময় মোজার উপর মাসেহ করেছেন। (সহীহ বুখারি ১৮২; সহীহ মুসলিম ২৭৪)
\end{quote}

\begin{itemize}
\item \textbf{সময়সীমা:} মুকিমের জন্য এক দিন-এক রাত (২৪ ঘণ্টা), মুসাফিরের জন্য তিন দিন-তিন রাত। (সহীহ মুসলিম ২৭৬)
\item মাসেহ শুরুর হিসাব ওযু ভাঙার পর প্রথম মাসেহ থেকে — এ নিয়ে মাযহাবগত সামান্য পার্থক্য আছে (আরও বিশদ: পার্ট ৫গ)।
\end{itemize}
\subsubsection*{২.২ গোসল (বড় অপবিত্রতা দূর করা)}

গোসল ফরজ হয়: \textbf{জুনুবি} (স্বপ্নদোষ/সহবাস/বীর্য নির্গমন) হলে, \textbf{হায়েয ও নিফাস} শেষে, এবং মৃত্যুর পর মাইয়েত গোসলের ক্ষেত্রে। (আরও বিশদ: পার্ট ৫গ)

\textbf{মূল আয়াত (সূরা আন-নিসা ৪:৪৩):}

\begin{center}\arab{يَـٰٓأَيُّهَا ٱلَّذِينَ ءَامَنُوا۟ لَا تَقْرَبُوا۟ ٱلصَّلَوٰةَ وَأَنتُمْ سُكَـٰرَىٰ حَتَّىٰ تَعْلَمُوا۟ مَا تَقُولُونَ وَلَا جُنُبًا إِلَّا عَابِرِى سَبِيلٍ حَتَّىٰ تَغْتَسِلُوا۟ ۚ وَإِن كُنتُم مَّرْضَىٰٓ أَوْ عَلَىٰ سَفَرٍ أَوْ جَآءَ أَحَدٌ مِّنكُم مِّنَ ٱلْغَآئِطِ أَوْ لَـٰمَسْتُمُ ٱلنِّسَآءَ فَلَمْ تَجِدُوا۟ مَآءً فَتَيَمَّمُوا۟ صَعِيدًا طَيِّبًا فَٱمْسَحُوا۟ بِوُجُوهِكُمْ وَأَيْدِيكُمْ ۗ إِنَّ ٱللَّهَ كَانَ عَفُوًّا غَفُورًا}\end{center}

\textbf{উচ্চারণ:} ইয়া- আয়্যুহাল্লাযীনা আ-মানূ, লা- তাকরাবুস-সলাতা ওয়া আনতুম সুকা-রা- হাত্তা- তা'লামূ মা- তাকূলূন, ওয়ালা- জুনুবান ইল্লা- 'আবিরী সাবীলিন হাত্তা- তাগতাছিলূ। ওয়া ইন কুনতুম মারদ্বা- আও 'আলা- সাফারিন, আও জাআ আহাদুম মিনকুম মিনাল-গাইত্বি, আও লা-মাসতুমুন-নিসাআ, ফালাম তাজিদূ মাআন, ফাতায়াম্মামূ স'য়ীদান ত্বাইয়িবান ফামসাহূ বিউজূহিকুম ওয়া আইদীকুম। ইন্নাল্লা-হা কা-না 'আফুওওয়ান গাফূরা।

\textbf{বাংলা অর্থ (\lat{pure}):}
\textbf{\lat{“}হে ঈমানদারগণ! তোমরা নেশাগ্রস্ত অবস্থায় নামাজের কাছে যেয়ো না — যতক্ষণ না বুঝতে পারো যা বলছ — এবং জুনুবি অবস্থাতেও (যেয়ো না), তবে পথ চলাকালীন (মুসাফির) হলে ভিন্ন — যতক্ষণ না গোসল করো। আর যদি অসুস্থ থাকো বা সফরে থাকো, অথবা কেউ পায়খানা থেকে আসে, অথবা স্ত্রীদের সাথে সহবাস করে থাকো, আর পানি না পাও, তবে পবিত্র মাটি দিয়ে তায়াম্মুম করো এবং তা দিয়ে মুখমণ্ডল ও দুই হাত মাসেহ করো। নিশ্চয়ই আল্লাহ ক্ষমাশীল, পরম ক্ষমাশীল।\lat{”}}

\textbf{গোসলের ফরজ (সর্বসম্মত মূল):} সমগ্র শরীরে পানি পৌঁছানো। শাফেয়ি ও হাম্বলি মতে \textbf{কুলি ও নাকে পানি} দেওয়াও ফরজ; হানাফি মতে তা সুন্নাত (পুরো শরীর ধোয়া ফরজ)। মালেকি মতে কুলি-নাকে পানি দেওয়া ফরজের অন্তর্ভুক্ত।

\textbf{সহীহ হাদিস (সহীহ বুখারি ২৪৮):}

\begin{quote}
আবু সালামাহ (রহ.) থেকে: আমি আয়েশা (রাঃ)-কে নবীজির (সা.) গোসলের পদ্ধতি জিজ্ঞেস করলাম। তিনি বললেন: নবীজি (সা.) প্রথমে বাম হাতে পানি নিয়ে (নাপাক) ধুয়ে ফেলতেন, তারপর জননাঙ্গ ধুয়ে নিতেন, তারপর ওযুর মতো ওযু করতেন, তারপর পুরো শরীরে পানি ঢালতেন।
\end{quote}

\subsubsection*{২.৩ তায়াম্মুম (পানি না পেলে / না পারলে)}

\textbf{ভিত্তি:} সূরা আন-নিসা ৪:৪৩ ও সূরা আল-মায়িদাহ ৫:৬ (উপরের দুটি আয়াত)।

\textbf{সহীহ হাদিস (সহীহ বুখারি ৩৪৮):}

\begin{quote}
ইমরান ইবনে হুসাইন (রাঃ) থেকে: নবীজি (সা.) এক ব্যক্তিকে জামাত থেকে আলাদা দাঁড়িয়ে থাকতে দেখে জিজ্ঞেস করলেন — নামাজে যোগ দাওনি কেন? সে বলল: আমি জুনুবি, আর পানি নেই। নবীজি (সা.) বললেন: \textbf{\lat{“}পবিত্র মাটি (স'য়ীদ) নিয়ে তায়াম্মুম করো — সেটাই তোমার জন্য যথেষ্ট।\lat{”}}
\end{quote}

\textbf{তায়াম্মুমের শর্ত ও পদ্ধতি (সংক্ষিপ্ত):}
\begin{enumerate}
\item পানি না পাওয়া বা অসুস্থতার কারণে পানি ব্যবহারে অক্ষম হওয়া (অথবা পানি পেতে বিপদ/বিলম্ব হওয়া)।
\item পবিত্র মাটি/মাটিজাতীয় বস্তুতে দুই হাত মেরে \textbf{চেহারায় মাসেহ}, তারপর \textbf{দুই হাতে} (কব্জি পর্যন্ত — মূল মত) মাসেহ করা।
\item \textbf{ধাক্কার সংখ্যা} (এক না দুই) ও \textbf{হাত পর্যন্ত সীমা} — এ নিয়ে মাযহাবগুলোর মধ্যে পার্থক্য আছে; মূল আমল একটাই: চেহারা ও দুই হাত মাসেহ।
\end{enumerate}
\begin{quote}
\textbf{নোট:} তায়াম্মুম \textbf{হাদাস দূর করে না}; এটি কেবল সেই ওয়াক্ত/আমলের জন্য পবিত্রতার অনুমতি। পানি পাওয়া গেলে আবার ওযু/গোসল আবশ্যক। (এটি সর্বসম্মত মূলনীতি — বিশদ: পার্ট ৫গ।)
\end{quote}

\subsubsection*{২.৪ শরীর, কাপড় ও স্থানের পবিত্রতা (নাজাস দূর করা)}

\begin{itemize}
\item কাপড়ে নাজাস (প্রস্রাব, মল, রক্ত, মদ ইত্যাদি) থাকলে তা \textbf{ধুয়ে} নামাজ পড়তে হবে। নববী যুগে মসজিদে কেউ প্রস্রাব করলে নবীজি (সা.) বলেছেন: \lat{“}...সেই জায়গায় এক বালতি পানি ঢেলে দাও\lat{”} (সহীহ বুখারি ২১৯, সহীহ মুসলিম ২৮৫)।
\item \textbf{পানি পবিত্র:} \lat{“}পানি পবিত্র, কোনো কিছুই তাকে নাপাক করে না\lat{”} (সুনানে আবু দাউদ ৬৬ — সহীহ/আলবানি)।
\item \textbf{ছোট পরিমাণ নাজাসের ক্ষমা} (মাফ-এ-মাগফুরাত): ফকিহদের মধ্যে আছে — যেমন রক্তের সামান্য ফোঁটা, হাঁটা-পথের কাদা-ময়লা ইত্যাদি। হানাফি মতে প্রস্রাব-পায়খানা ব্যতীত অন্যান্য নাজিসের পরিমাণ \textbf{দিরহামের চেয়ে কম} হলে ক্ষমা (নামাজের জন্য মাপসামা করা হয়)। শাফেয়ি মতে কোনো নাজিসই 'ক্ষমার' নির্দিষ্ট পরিমাণের বর্ণনা নির্ভরযোগ্য নয় — বরং সাধারণত ধোয়া বাধ্যতামূলক; তবে কষ্টকর অবস্থা থেকে রেহাই দেওয়া হয়েছে। \textbf{এই পার্থক্যটি বাস্তব এবং বৈধ।}
\item \textbf{শরীর/স্থানে নাজাস:} একই নিয়ম — ধুয়ে পরিষ্কার করা।
\end{itemize}
\medskip\hrule\medskip

\subsection*{৩. সময় (ওয়াক্ত) — নামাজের দ্বিতীয় মূল শর্ত}

\textbf{মূল আয়াত (সূরা আন-নিসা ৪:১০৩):}

\begin{center}\arab{فَإِذَا قَضَيْتُمُ ٱلصَّلَوٰةَ فَٱذْكُرُوا۟ ٱللَّهَ قِيَـٰمًا وَقُعُودًا وَعَلَىٰ جُنُوبِكُمْ ۚ فَإِذَا ٱطْمَأْنَنتُمْ فَأَقِيمُوا۟ ٱلصَّلَوٰةَ ۚ إِنَّ ٱلصَّلَوٰةَ كَانَتْ عَلَى ٱلْمُؤْمِنِينَ كِتَـٰبًا مَّوْقُوتًا}\end{center}

\textbf{উচ্চারণ:} ফাইযা- কাজাইতুমুস-সলাতা ফাযকুরুল্লা-হা ক্বিয়ামাও ওয়া ক্বু'ঊদাও ওয়া 'আলা- জুনূবিকুম। ফাইযা-ত্বমা'নান্তুম ফাআক্বীমুস-সলাতা। ইন্নাস-সলাতা কা-নাত 'আলাল-মু'মিনীনা কিতা-বাম মাওকূতা।

\textbf{বাংলা অর্থ (\lat{pure}):}
\textbf{\lat{“}অতঃপর যখন তোমরা নামাজ শেষ করো, তখন আল্লাহকে স্মরণ করো দাঁড়ানো, বসা ও কাত হয়ে শুয়ে অবস্থায়। আর যখন তোমরা নিরাপদ হবে, তখন নামাজ (পূর্ণ) কায়েম করো। নিশ্চয়ই নামাজ নির্দিষ্ট সময়ে মুমিনদের উপর ফরজ করা হয়েছে।\lat{”}}

\textbf{সহীহ হাদিস — নামাজের সময়ের সীমা (সুনানে তিরমিজি ১৪৯ — হাসান):}

\begin{quote}
ইবনে আব্বাস (রাঃ) থেকে: নবীজি (সা.) বলেছেন — \textbf{\lat{“}জিবরিল (আঃ) আমার সাথে বাইতুল্লাহর কাছে দুইবার নামাজ পড়েছেন। প্রথমবার যোহর পড়লেন যখন ছায়া জুতার ফিতার সমান হলো; আসর পড়লেন যখন প্রতিটি বস্তুর ছায়া নিজের সমান হলো; মাগরিব পড়লেন যখন সূর্য অস্ত গেল (রোজাদার ইফতার করে); ইশা পড়লেন যখন শাফাক (লাল আভা) বিদায় নিল; ফজর পড়লেন যখন ফজর উদিত হলো (রোজাদারের খাওয়া হারাম হলো)। দ্বিতীয়বার যোহর পড়লেন যখন প্রতিটি বস্তুর ছায়া নিজের সমান হলো; আসর পড়লেন যখন ছায়া নিজের দ্বিগুণ হলো; মাগরিব প্রথমবারের মতোই সময়ে; ইশা পড়লেন যখন রাতের এক-তৃতীয়াংশ গেল; ফজর পড়লেন যখন (আকাশ) উজ্জ্বল হয়ে উঠল। তারপর জিবরিল বললেন: হে মুহাম্মাদ! এগুলো তোমার পূর্ববর্তী নবীদের সময়; আর (প্রকৃত) সময় হলো এই দুই সময়ের মাঝখানে।\lat{”}}
\end{quote}

\textbf{ওয়াক্তের সারণি (সংক্ষিপ্ত):}

\begin{table}[h]\centering\small
\begin{tabular}{lll}
নামাজ & শুরুর সময় & শেষ সময় \\
\hline
ফজর & সুবহে সাদিক (ভোরের সাদা আলো) & সূর্যোদয় \\
যোহর & সূর্য পশ্চিমে ঢলে পড়া & জমহুর (শাফেয়ি-মালেকি-হাম্বলি): ছায়া নিজের সমান; হানাফি: ছায়া নিজের দ্বিগুণ \\
আসর & জমহুর: ছায়া নিজের সমান; হানাফি: ছায়া দ্বিগুণ & সূর্য হলদে/লাল হয়ে অস্ত যাওয়া \\
মাগরিব & সূর্যাস্ত & শাফাক (লাল আভা) অদৃশ্য হওয়া \\
ইশা & শাফাক অদৃশ্য হওয়া & মধ্যরাত (মাযহাবগত পার্থক্য) \\
\hline
\end{tabular}
\end{table}

\begin{quote}
\textbf{যোহর ও আসরের সন্ধিক্ষণে বাস্তব ইখতিলাফ:} শাফেয়ি-মালেকি-হাম্বলি মতে আসরের শুরু যখন ছায়া নিজের সমান হয়; হানাফি মতে যখন দ্বিগুণ হয়। ফলে \textbf{হানাফি মাযহাবে আসরের ওয়াক্ত পরে শুরু হয়}। এটা যোহর/আসরের সময়সূচিতে কার্যত পার্থক্য সৃষ্টি করে — বৈধ ইখতিলাফ, কোনো পক্ষ 'ভুল' নয়।
\end{quote}

\textbf{আরও সহীহ হাদিস — সময়ের ফজিলত (সহীহ বুখারি ৫২৭):}

\begin{quote}
আবদুল্লাহ ইবনে মাসউদ (রাঃ) থেকে: আমি নবীজি (সা.)-কে জিজ্ঞেস করলাম — কোন আমল আল্লাহর কাছে সবচেয়ে প্রিয়? তিনি বললেন: \textbf{\lat{“}সময়মতো (প্রথম ওয়াক্তে) নামাজ।\lat{”}} বললাম: এরপর? বললেন: \lat{“}পিতামাতার প্রতি সদ্ব্যবহার।\lat{”} বললাম: এরপর? বললেন: \lat{“}আল্লাহর পথে জিহাদ।\lat{”}
\end{quote}

\begin{quote}
\textbf{মাঝের নামাজ (আসর) সংরক্ষণ:} \lat{“}তোমরা সব নামাজের হেফাজত করো, বিশেষ করে মধ্যবর্তী নামাজের\lat{”} (সূরা আল-বাকারাহ ২:২৩৮)। তাফসিরকারদের মতে 'মধ্যমা নামাজ' = আসর (প্রসিদ্ধ মত)। নামাজ নির্দিষ্ট সময়ে আদায়ে অলসতা হারাম-সদৃশ (মাযহাবগত দৃষ্টিভঙ্গি পার্ট ১-এ)।
\end{quote}

\textbf{ব্যবহারিক নিয়ম:}
\begin{itemize}
\item প্রথম ওয়াক্তে নামাজ আদায় করা সর্বোত্তম (বুখারি ৫২৭; উত্তাপে যোহর দেরি করার সুন্নাতের ব্যতিক্রম — পার্ট ৮)।
\item ঘুমের কারণে বা ভুলে ফরজ ছুটে গেলে জাগার/মনে হওয়ার সঙ্গে সঙ্গে \textbf{কাযা} আদায় করতে হবে (সহীহ বুখারি ৫৯৭, সহীহ মুসলিম ৬৮৪)।
\end{itemize}
\medskip\hrule\medskip

\subsection*{৪. কিবলা — তৃতীয় মূল শর্ত}

\textbf{মূল আয়াত (সূরা আল-বাকারাহ ২:১৪৪):}

\begin{center}\arab{قَدْ نَرَىٰ تَقَلُّبَ وَجْهِكَ فِى ٱلسَّمَآءِ ۖ فَلَنُوَلِّيَنَّكَ قِبْلَةً تَرْضَىٰهَا ۚ فَوَلِّ وَجْهَكَ شَطْرَ ٱلْمَسْجِدِ ٱلْحَرَامِ ۚ وَحَيْثُ مَا كُنتُمْ فَوَلُّوا۟ وُجُوهَكُمْ شَطْرَهُۥ ۗ وَإِنَّ ٱلَّذِينَ أُوتُوا۟ ٱلْكِتَـٰبَ لَيَعْلَمُونَ أَنَّهُ ٱلْحَقُّ مِن رَّبِّهِمْ ۗ وَمَا ٱللَّهُ بِغَـٰفِلٍ عَمَّا يَعْمَلُونَ}\end{center}

\textbf{উচ্চারণ:} ক্বাদ নারা- তাক্বাল্লুবা ওয়াজহিকা ফিস-সামা-য়ি, ফালানুওয়াল্লিয়ান্নাকা ক্বিবলাতান তারদ্বা-হা-। ফাওয়াল্লি ওয়াজহাকা শাত্বরাল-মাসজিদিল-হারাম, ওয়া হাইসু মা- কুনতুম ফাওয়াল্লূ উজূহাকুম শাত্বরাহূ। ওয়া ইন্নাল্লাযীনা উতুল-কিতা-বা লাইয়া'লামূনা আন্নাহুল-হাক্কু মির-রাব্বিহিম। ওয়া মাল্লা-হু বিগা-ফিলিন 'আম্মা- ইয়া'মালূন।

\textbf{বাংলা অর্থ (\lat{pure}):}
\textbf{\lat{“}আমি তো দেখছি তোমার মুখ আকাশের দিকে ফেরানো হচ্ছে। আমি অবশ্যই তোমাকে এমন কিবলার দিকে ফেরাব যা তুমি পছন্দ করবে। অতএব তোমার মুখ মসজিদুল হারামের দিকে ফেরাও। আর যেখানেই তোমরা থাকো, তোমরা সবাই তোমাদের মুখ সেদিকে ফেরাও। নিশ্চয়ই যাদের কিতাব দেওয়া হয়েছে তারা জানে এটি তাদের রবের পক্ষ থেকে সত্য। আর আল্লাহ তাদের আমল সম্পর্কে গাফিল নন।\lat{”}}

\textbf{ব্যাখ্যা:} কিবলা পরিবর্তনের ঘটনা — নবীজি (সা.) ষোলো মাস বাইতুল-মাকদিসের দিকে ফিরে নামাজ পড়ার পর কাবার দিকে ফেরার নির্দেশ এসেছে। ফরজ নামাজে \textbf{কিবলামুখী হওয়া আবশ্যক}।

\textbf{সহীহ হাদিস (সহীহ বুখারি ৪০০):}

\begin{quote}
জাবির (রাঃ) থেকে: নবীজি (সা.) সওয়ারিতে (বাহনে) \textbf{নফল} নামাজ পড়তেন যেদিকে বাহন মুখ করে থাকত; কিন্তু \textbf{ফরজ} নামাজের ইচ্ছা করলে নেমে কিবলামুখী হয়ে পড়তেন।
\end{quote}

\begin{quote}
\textbf{ব্যবহারিক নিয়ম:} যেখানে যেখানেই থাকুন, যথাসাধ্য কিবলার দিক অনুমান করে ফরজ নামাজ কিবলামুখী হতেই হবে। কিবলা নিয়ে সন্দেহ হলে ইজতিহাদ (চেষ্টা/অনুমান) করা বৈধ। অসুস্থতা/ভয়-ভীতি ও সফরে নির্দিষ্ট অবস্থায় শিথিলতা আছে (পার্ট ৫গ)।
\end{quote}

\medskip\hrule\medskip

\subsection*{৫. সতর (আওরাত) ও পোশাক}

\textbf{মূল আয়াত (সূরা আল-আ'রাফ ৭:৩১):}

\begin{center}\arab{۞ يَـٰبَنِىٓ ءَادَمَ خُذُوا۟ زِينَتَكُمْ عِندَ كُلِّ مَسْجِدٍ وَكُلُوا۟ وَٱشْرَبُوا۟ وَلَا تُسْرِفُوٓا۟ ۚ إِنَّهُۥ لَا يُحِبُّ ٱلْمُسْرِفِينَ}\end{center}

\textbf{উচ্চারণ:} ইয়া- বানী- আদামা খুযূ যীনাতাকুম 'ইন্দা কুল্লি মাসজিদিন, ওয়া কুলূ ওয়াশরাবূ ওয়ালা- তুসরিফূ। ইন্নাহূ লা- ইউহিব্বুল-মুসরিফীন।

\textbf{বাংলা অর্থ (\lat{pure}):}
\textbf{\lat{“}হে আদম-সন্তান! প্রত্যেক মসজিদে (নামাজের সময়) তোমরা তোমাদের সৌন্দর্য (পরিচ্ছন্ন-শালীন পোশাক) গ্রহণ করো। আর খাও, পান করো, কিন্তু অপচয় কোরো না। নিশ্চয়ই তিনি অপচয়কারীদের পছন্দ করেন না।\lat{”}}

\textbf{সহীহ হাদিস — পুরুষের জন্য (সহীহ বুখারি ৩৫৮):}

\begin{quote}
আবু হুরাইরা (রাঃ) থেকে: এক ব্যক্তি নবীজি (সা.)-কে এক কাপড়ে নামাজ পড়ার কথা জিজ্ঞেস করল। নবীজি (সা.) বললেন: \textbf{\lat{“}তোমাদের সবার কি দুটো করে কাপড় আছে?\lat{”}} (অর্থাৎ — নামাজে অন্তত যথাযথ সতর ঢাকার মতো ব্যবস্থা করতে হবে; আয়েশা (রাঃ) থেকে আরেক বর্ণনায় — ঠান্ডায় নবীজি (সা.) একই কাপড়ে দুই দিকে জড়িয়ে নামাজ পড়েছেন।)
\end{quote}

\textbf{সহীহ হাদিস — নারীর জন্য (সুনানে আবু দাউদ ৬৪১ — সহীহ/আলবানি):}

\begin{quote}
আয়েশা (রাঃ) থেকে: নবীজি (সা.) বলেছেন — \textbf{\lat{“}বয়োঃপ্রাপ্তা নারীর নামাজ খিমার (মাথা-গলা ঢাকা কাপড়) ছাড়া কবুল হয় না।\lat{”}}
\end{quote}

\textbf{সতরের পরিধি — মাযহাবগত অবস্থান (সংক্ষেপ):}

\begin{table}[h]\centering\small
\begin{tabular}{ll}
ব্যক্তি & আওরাত (নামাজে ঢাকতে হবে) \\
\hline
পুরুষ & নাভি থেকে হাঁটু পর্যন্ত (মাযহাবগুলোর প্রসিদ্ধ মত — তবে সরাসরি এই সীমা নির্ধারণী হাদিস দুর্বল; উভয় প্রান্তে সামান্য ছাড় দেওয়া হয়)। কাঁধ ঢাকা উত্তম (বুখারি ৩৫৮-এর আলোকে) \\
নারী & মুখ ও দুই হাত ছাড়া সমগ্র শরীর; হাম্বলি ও হানাফি মতেও পায়ের পাতা ঢাকা উত্তম; মালেকি মতে চেহারা ছাড়া সব (পায়ের নিচে শিথিলতা আছে) \\
\hline
\end{tabular}
\end{table}

\begin{quote}
\textbf{ভিত্তি:} সূরা আল-আ'রাফ ৭:৩১, সূরা আন-নূর ২৪:৩১, এবং উপরের সহীহ হাদিস। 'নাভি-হাঁটু' সীমা সম্পর্কিত হাদিসটি দুর্বল — তাই এটিকে 'সব মাযহাবের গ্রহণযোগ্য বাস্তব নিয়ম' হিসেবে তুলে ধরা হয়, সরাসরি কুরআনের নিষ্ক্রিয় শব্দ হিসেবে নয়।
\end{quote}

\textbf{ব্যবহারিক নিয়ম:}
\begin{itemize}
\item নামাজে ঢিলেঢালা, সতর ঢাকে এমন পোশাক; স্বচ্ছ/আঁটসাঁট পোশাকে সতরের হুকুম নিয়ে মতভেদ আছে (শরীর ঢাকা মনে হলে বৈধ — তবে সুন্নাত-সম্মত নয়)।
\item মসজিদে যাওয়ার আগে পোশাক-পরিচ্ছন্নতার যত্ন নেওয়া সুন্নাত (সূরা আল-আ'রাফ ৭:৩১)।
\end{itemize}
\medskip\hrule\medskip

\subsection*{৬. আযান ও ইকামত}

\begin{itemize}
\item \textbf{আযান} পুরুষদের ওয়াক্তভিত্তিক জামাতের জন্য ঘোষণা; ফরজ নামাজের সময় নির্দেশ।
\item \textbf{সহীহ হাদিস (সহীহ বুখারি ৬১১):} \lat{“}যখন তোমরা আযান শুনো, মুয়াযযিন যা বলে তোমরাও তাই বলো।\lat{”}
\item \textbf{সহীহ হাদিস (সহীহ বুখারি ৬১৫):} \lat{“}লোকেরা যদি জানত আযান ও প্রথম কাতারে কী (ফজিলত) আছে — আর তা পেতে লটারি ছাড়া উপায় না পেত, তবে তারা লটারি করত। আর যদি জানত (যোহরের) প্রথম ওয়াক্তে ও (ইশা-ফজরের) জামাতে কী আছে, তবে হামাগুড়ি দিয়েও এসে পড়ত।\lat{”}
\item \textbf{আযানের শব্দাবলি (সুনানে আবু দাউদ ৪৯৯ — হাসান সহীহ/আলবানি):} আবদুল্লাহ ইবনে যায়েদ (রাঃ)-এর স্বপ্নে শেখানো আযান — \lat{“}আল্লাহু আকবার (৪ বার), আশহাদু আল্লা ইলাহা ইল্লাল্লাহ (২ বার), আশহাদু আন্না মুহাম্মাদার রাসূলুল্লাহ (২ বার), হাইয়া 'আলাস সালাহ (২ বার), হাইয়া 'আলাল ফালাহ (২ বার), আল্লাহু আকবার (২ বার), লা ইলাহা ইল্লাল্লাহ (১ বার)।\lat{”} ইকামতে \lat{“}হাইয়া 'আলাল ফালাহ\lat{”}-এর পর \lat{“}ক্বাদ ক্বামাতিস সালাহ\lat{”} ২ বার। ফজরের আযানে \lat{“}হাইয়া 'আলাল ফালাহ\lat{”}-এর পর \lat{“}আসসালাতু খাইরুম মিনান নাউম\lat{”} ২ বার। (এই শব্দটি ফজরের বিশেষ সুন্নাত।)
\item \textbf{মুআযযিনের জবাব:} \lat{“}হাইয়া 'আলাস সালাহ\lat{”} ও \lat{“}হাইয়া 'আলাল ফালাহ\lat{”} শুনে \lat{“}লা হাওলা ওয়ালা কু\lat{ww}আতা ইল্লা বিল্লাহ\lat{”} বলা; \lat{“}আসসালাতু খাইরুম মিনান নাউম\lat{”}-এ \lat{“}সাদাকত্বা ওয়া বারারতা\lat{”} বলা — এসব সুন্নাত (সহীহ মুসলিম ৩৮৫)।
\end{itemize}
\begin{quote}
\textbf{মাযহাবগত পার্থক্য:} আযান-ইকামত \textbf{পুরুষদের} জন্য সুন্নাতে মুআক্কাদা (ফরজ কেফায়া-সমান); নারীদের জন্য নয় — এ নিয়ে প্রায় ঐকমত্য; মালেকি মতে নারীর আযান মাকরূহ। (বিশদ: পার্ট ৫গ-এ নারী-বিষয়ক আলোচনা।)
\end{quote}

\medskip\hrule\medskip

\subsection*{৭. নিয়ত (\lat{Intention})}

\textbf{মূল হাদিস (সহীহ বুখারি ১; সহীহ মুসলিম ১৯০৭):}

\begin{quote}
উমার ইবনুল খাত্তাব (রাঃ) থেকে: রাসূলুল্লাহ (সা.) বলেছেন — \textbf{\lat{“}প্রতিটি আমল নিয়তের উপর নির্ভরশীল। প্রত্যেক ব্যক্তি যা নিয়ত করেছে তা-ই পাবে। যে দুনিয়া লাভের জন্য কিংবা বিবাহের জন্য হিজরত করেছে, তার হিজরত সে জন্যই।\lat{”}}
\end{quote}

\textbf{ব্যবহারিক নিয়ম:}
\begin{itemize}
\item নিয়তের স্থান \textbf{হৃদয়}। মুখে নিয়ত বলা (লাফযিয়্যাহ) সম্পর্কে আলেমদের মধ্যে মতভেদ আছে: প্রসিদ্ধ ও নির্ভরযোগ্য মত হলো — \textbf{মুখে বলা জরুরি নয়}; বরং কেউ কেউ এটাকে বিদ'আত-নিকট বলেছেন (সালাফের আমলে নিয়ত মুখে বলা হতো না)। তবে অনেক মাযহাব-গ্রন্থে মুখে নিয়ত বলা মুস্তাহাব বলা হয়েছে। \textbf{নিরাপদ ও সুন্নাহ-সম্মত:} হৃদয়ে নিয়ত + তাকবির দিয়ে নামাজ শুরু।
\item নিয়তের সাথে \textbf{কোন নামাজ} (যেমন যোহর/ফরজ/কাযা) তা নির্ধারণ করা জরুরি; ইমাম হলে ইমামতি ও মুক্তাদি হলে ইকতিদা — হৃদয়ে নির্ধারণ করা ভালো।
\end{itemize}
\medskip\hrule\medskip

\subsection*{৮. স্থান (সাল্লার জায়গা)}

\begin{itemize}
\item \textbf{সহীহ হাদিস (সহীহ বুখারি ৩৩৫):} \lat{“}পৃথিবীকে আমার জন্য মসজিদ ও পবিত্রতার মাধ্যম বানানো হয়েছে। সুতরাং আমার উম্মতের যে কেউ যেকোনো স্থানে নামাজের সময় হলে সেখানে নামাজ পড়বে।\lat{”} — এর মানে \textbf{পবিত্র ও নিরাপদ যেকোনো জায়গা} মসজিদ হিসেবে গণ্য; তবে ভেজা-নাপাক জায়গায় সিজদা করা যাবে না।
\item \textbf{মসজিদে প্রবেশের আদব:} মসজিদে ঢুকলে দুই রাকাত তাহিয়্যাতুল মসজিদ (সহীহ বুখারি ৪৪৪; সহীহ মুসলিম ৭১৪)।
\item \textbf{সুতরা (সামনে কিছু রাখা):} নামাজীর সামনে সুতরা রাখা ও না রাখা — পার্ট ৭ক-এ বিস্তারিত।
\end{itemize}
\medskip\hrule\medskip

\subsection*{সারসংক্ষেপ — এক নজরে}

\begin{table}[h]\centering\small
\begin{tabular}{lll}
শর্ত & মূল দলিল & সতর্কতা \\
\hline
পবিত্রতা (ওযু/গোসল/তায়াম্মুম) & ৫:৬, ৪:৪৩; বুখারি ১৩৫, ২৪৮, ৩৪৮ & ওযু ভঙ্গের তালিকা মাযহাবভেদে ভিন্ন \\
সময় & ৪:১০৩; তিরমিজি ১৪৯ (হাসান) & আসরের শুরুতে হানাফি-জমহুর পার্থক্য \\
কিবলা & ২:১৪৪; বুখারি ৪০০ & সন্দেহে ইজতিহাদ বৈধ \\
সতর & ৭:৩১; বুখারি ৩৫৮; আবু দাউদ ৬৪১ & 'নাভি-হাঁটু' হাদিস দুর্বল — বাস্তব নিয়ম হিসেবে বহাল \\
আযান & বুখারি ৬১১, ৬১৫; আবু দাউদ ৪৯৯ & পুরুষের জন্য সুন্নাতে মুআক্কাদা \\
নিয়ত & বুখারি ১ & হৃদয়ে; মুখে বলা বিতর্কিত \\
স্থান & বুখারি ৩৩৫ & পবিত্র স্থান প্রয়োজন \\
\hline
\end{tabular}
\end{table}

\medskip\hrule\medskip

\subsection*{উৎসগ্রন্থ তালিকা (এই পার্টে ব্যবহৃত)}

\begin{itemize}
\item কুরআন: সূরা আল-মায়িদাহ ৫:৬; আন-নিসা ৪:৪৩, ৪:১০৩; আল-বাকারাহ ২:১৪৪, ২:২৩৮; আল-আ'রাফ ৭:৩১; আল-মুদাসসির ৭৪:৪-৫; আন-নূর ২৪:৩১
\item সহীহ বুখারি: ১, ১৩৫, ১৮২, ২১৯, ২৪৮, ৩৩৫, ৩৪৮, ৩৫৮, ৪০০, ৫২৭, ৫৯৭, ৬১১, ৬১৫, ৬৩১, ৮৫৮
\item সহীহ মুসলিম: ২৭৪, ২৭৬, ২৮৫, ৩৬০, ৩৮৫, ৫৫৩, ৬৬৭, ৬৮৪, ১৯০৭
\item সুনানে আবু দাউদ: ৬৬, ১৮১, ৪৯৯, ৬৪১
\item সুনানে তিরমিজি: ৮২, ১৪৯, ১৪২৩
\item সিফাতু সালাতিন-নবী (আলবানি) — পদ্ধতি ও হাদিস শনাক্তকরণের রেফারেন্স
\end{itemize}
\begin{quote}
\textbf{দ্রষ্টব্য:} যেসব বিষয় এখানে সংক্ষিপ্ত রাখা হয়েছে — মাসেহ খুফফাইনের বিস্তার, তায়াম্মুমের পূর্ণ শর্ত, হায়েয-নিফাস, নারী-জামাত, মুসাফিরের বিধান — সেগুলোর পূর্ণাঙ্গ আলোচনা পার্ট ৫গ ও ৫ঘ-এ।
\end{quote}


\end{document}
